\documentclass[aspectratio=169,11pt]{beamer}

%==============================================================================
% Theme
%==============================================================================

\usetheme{Madrid}
\usecolortheme{default}
\setbeamercolor{frametitle}{fg=white,bg=polimi}
\setbeamertemplate{navigation symbols}{}
\setbeamertemplate{footline}[frame number]




%==============================================================================
% Packages
%==============================================================================

\usepackage[utf8]{inputenc}
\usepackage[T1]{fontenc}
\usepackage{graphicx}
\usepackage{booktabs}
\usepackage{amsmath}
\usepackage{multicol}
\usepackage{tikz}
\usepackage{xcolor}
\usepackage{hyperref}

%==============================================================================
% Image handling
%==============================================================================

% Images are searched in the Images folder.
\graphicspath{{Images/}}

% Safe image command:
% If the image exists, it is displayed.
% If it does not exist, a placeholder is displayed instead.
%
% This prevents missing-image errors from stopping compilation.
\newcommand{\placeholderimage}[2][0.8\textwidth]{%
    \IfFileExists{#2}{%
        \includegraphics[width=#1]{#2}%
    }{%
        \fbox{%
            \begin{minipage}[c][0.22\textheight][c]{#1}
                \centering
                \texttt{\detokenize{#2}}\\[0.3cm]
                \footnotesize Image placeholder
            \end{minipage}%
        }%
    }%
}

%==============================================================================
% Colors
%==============================================================================

\definecolor{polimi}{RGB}{0,45,100}
\definecolor{fcosblue}{RGB}{33,150,243}
\definecolor{chestred}{RGB}{180,60,60}
\definecolor{imagegreen}{RGB}{70,140,80}

% General Beamer colors
\setbeamercolor{structure}{fg=polimi}

% Frame titles: WHITE text on POLIMI BLUE background
\setbeamercolor{frametitle}{
    fg=white,
    bg=polimi
}

% Title page
\setbeamercolor{title}{
    fg=white,
    bg=polimi
}
%==============================================================================
% Title
%==============================================================================

\title{Pneumonia Detection Using\\
Anchor-Free Object Detection}

\subtitle{An FCOS-based study with ResNet backbones}

\author[Francesco Rausa]{Francesco Rausa \\ 10833982}\\

\institute[PoliMI]{
    Numerical Analysis for Machine Learning\\
    Politecnico di Milano
}

\date{2025--2026}

% If you have the Politecnico logo, you can uncomment this line:
% \titlegraphic{\includegraphics[height=1.4cm]{logo_polimi.png}}

%==============================================================================
% Document
%==============================================================================

\begin{document}

%==============================================================================
% TITLE
%==============================================================================

\begin{frame}
    \titlepage
\end{frame}

%------------------------------------------------------------------------------

\begin{frame}{Outline}
    \tableofcontents
\end{frame}

%==============================================================================
% 1. INTRODUCTION
%==============================================================================

\section{Introduction}

%------------------------------------------------------------------------------

\begin{frame}{Motivation and Objective}

\begin{columns}

\begin{column}{0.56\textwidth}

\begin{itemize}

    \item Pneumonia detection from chest X-ray images is a challenging
    medical imaging task.

    \item The objective is not only to determine whether pneumonia is present,
    but also to \textbf{localize the affected region}.

    \item Pneumonia findings can vary substantially in
    \textbf{size, shape and position}.

    \item This motivates the use of an object detection framework rather than
    image-level classification alone.

\end{itemize}

\vspace{0.3cm}

\textbf{Main question:}

\begin{quote}
How do backbone depth and pretraining affect an
anchor-free pneumonia detector?
\end{quote}

\end{column}

\begin{column}{0.40\textwidth}

\centering

\placeholderimage[\textwidth]{Images/chest_xray_example.jpg}

\vspace{0.2cm}

{\footnotesize Example chest X-ray with pneumonia localization}

\end{column}

\end{columns}

\end{frame}

\begin{frame}{From Anchor-Based to Anchor-Free Detection}

\begin{columns}[t]

%------------------------------------------------------------------------------
% Anchor-based
%------------------------------------------------------------------------------

\begin{column}{0.47\textwidth}

\begin{block}{Anchor-Based}
\footnotesize
\begin{itemize}
    \setlength\itemsep{2pt}
    \item Predefined bounding boxes
    \item Classify anchors and regress box offsets
    \item Requires anchor sizes and aspect ratios
\end{itemize}
\end{block}

\end{column}

%------------------------------------------------------------------------------
% Anchor-free
%------------------------------------------------------------------------------

\begin{column}{0.47\textwidth}

\begin{block}{Anchor-Free}
\footnotesize
\begin{itemize}
    \setlength\itemsep{2pt}
    \item No predefined bounding boxes
    \item Predictions from spatial locations
    \item FCOS: class + $(l,t,r,b)$ + centerness
\end{itemize}
\end{block}

\end{column}

\end{columns}

\vspace{0.05cm}

\begin{center}
    \placeholderimage[0.45\textwidth]{Images/free_vs_anchor.png}
\end{center}

\vspace{0.02cm}

\centering
\textbf{This project focuses on the anchor-free formulation.}

\end{frame}
%------------------------------------------------------------------------------

\begin{frame}{From the Reference Study to the Project Implementation}

\begin{columns}

\begin{column}{0.46\textwidth}

\begin{block}{Reference}
Wu et al.~(2024)
\end{block}

\vspace{0.3cm}

\begin{itemize}
    \item Pneumonia detection on the RSNA dataset
    \item Anchor-free detection framework
    \item Feature pyramid
    \item Center and scale prediction
\end{itemize}

\end{column}

\begin{column}{0.46\textwidth}

\begin{block}{This project}
FCOS-based implementation
\end{block}

\vspace{0.3cm}

\begin{itemize}
    \item ResNet-50 / ResNet-101
    \item Feature Pyramid Network
    \item Classification
    \item LTRB regression
    \item Centerness
\end{itemize}

\end{column}

\end{columns}

\vspace{0.5cm}

\centering
\textbf{
Goal: study the FCOS formulation itself rather than reproduce the
reference implementation as a black box.
}

\end{frame}


%==============================================================================
% 2. ARCHITECTURE
%==============================================================================

\section{Architecture}

%------------------------------------------------------------------------------

\begin{frame}{Overall Architecture}

\begin{center}

\placeholderimage[0.70\textwidth]{Images/overall_architecture.png}

\end{center}

\vspace{0.2cm}

\centering

\textbf{Input}
$\rightarrow$
\textbf{ResNet}
$\rightarrow$
\textbf{FPN}
$\rightarrow$
\textbf{Detection Heads}

\end{frame}

%------------------------------------------------------------------------------

%==============================================================================
% RESNET
%==============================================================================

\section{ResNet Backbone}

%------------------------------------------------------------------------------

\begin{frame}{Why ResNet?}

\begin{columns}

\begin{column}{0.48\textwidth}

\begin{itemize}
    \item Increasing depth does not always improve training.
    
    \item He et al. observed a \textbf{degradation problem}:
    deeper plain networks can have \textbf{higher training error}.
    
    \item The effect is not simply due to overfitting.
    
    \item Main issue: \textbf{optimization difficulty}.
\end{itemize}

\end{column}

\begin{column}{0.48\textwidth}

\centering

\placeholderimage[0.55\textwidth]{Images/resnet_degradation_plain.png}

\vspace{0.15cm}

\placeholderimage[0.55\textwidth]{Images/resnet_degradation_resnet.png}

\end{column}

\end{columns}

\vspace{0.15cm}

\begin{block}{Key idea}
ResNet introduces \textbf{residual learning} and shortcut connections
to make deeper networks easier to optimize.
\end{block}

\end{frame}


%------------------------------------------------------------------------------

\begin{frame}{How Residual Learning Works}

\begin{columns}

\begin{column}{0.48\textwidth}

\begin{block}{Plain formulation}

\[
H(x)
\]

The stacked layers directly learn the desired transformation.

\end{block}

\vspace{0.25cm}

\begin{block}{Residual formulation}

\[
F(x) = H(x) - x
\]

\centering so that

\[
H(x) = F(x) + x
\]

\end{block}

\end{column}

\begin{column}{0.48\textwidth}

\centering

\placeholderimage[\textwidth]{Images/residual_block.png}

\vspace{0.2cm}

{\footnotesize
Shortcut connection + residual branch
}

\end{column}

\end{columns}

\vspace{0.25cm}

\centering

\textbf{The shortcut provides a direct identity path through the network.}

\end{frame}


%------------------------------------------------------------------------------

\begin{frame}{The Bottleneck Residual Block}

\begin{columns}

\begin{column}{0.42\textwidth}

\centering

\placeholderimage[\textwidth]{Images/bottleneck_block.png}

\end{column}

\begin{column}{0.53\textwidth}

\begin{block}{Three-layer residual branch}

\begin{enumerate}
    \item $1\times1$ convolution:
    \textbf{reduce} channels
    
    \item $3\times3$ convolution:
    spatial feature extraction
    
    \item $1\times1$ convolution:
    \textbf{restore} channels
\end{enumerate}

\end{block}

\vspace{0.25cm}

\[
C_{\mathrm{in}}
\;\rightarrow\;
C_{\mathrm{bottleneck}}
\;\rightarrow\;
C_{\mathrm{out}}
\]

\vspace{0.2cm}

\centering
\textbf{Why? Efficient deeper networks}

\end{column}

\end{columns}

\vspace{0.25cm}

{\footnotesize
For ResNet-50/101, the bottleneck uses $1\times1$, $3\times3$, and
$1\times1$ convolutions.
}

\end{frame}



%------------------------------------------------------------------------------

\begin{frame}{ResNet-50 and ResNet-101}

\begin{columns}[t]

\begin{column}{0.68\textwidth}

\centering

\begin{tabular}{lccccc}
\toprule
\textbf{Stage} & \textbf{R50} & \textbf{R101} &
\textbf{Resolution} & \textbf{Channels} \\
\midrule
Input       & -- & -- & $512\times512$ & 3 \\
Conv1       & -- & -- & $256\times256$ & 64 \\
MaxPool     & -- & -- & $128\times128$ & 64 \\
$conv2_x$   & 3  & 3  & $128\times128$ & 256 \\
$conv3_x$   & 4  & 4  & $64\times64$   & 512 \\
$conv4_x$   & 6  & 23 & $32\times32$   & 1024 \\
$conv5_x$   & 3  & 3  & $16\times16$   & 2048 \\
\bottomrule
\end{tabular}

\vspace{0.25cm}

{\footnotesize
Conv1: $7\times7$, 64 channels, stride 2
\qquad
MaxPool: $3\times3$, stride 2
}

\end{column}

\begin{column}{0.27\textwidth}

\begin{block}{Key idea}
\footnotesize
\begin{itemize}
    \item Same bottleneck architecture
    \item R101 is deeper mainly in $conv4_x$
    \item Resolution decreases
    \item Channels increase
\end{itemize}
\end{block}

\vspace{0.15cm}

\[
C_3,\;C_4,\;C_5
\]

\footnotesize
Feature maps extracted from the backbone

\end{column}

\end{columns}

\end{frame}





%------------------------------------------------------------------------------

\begin{frame}{Why Multi-Scale Features?}

\begin{columns}[t]

\begin{column}{0.47\textwidth}

\begin{block}{Deep layers}
\begin{itemize}
    \item High-level semantic information
    \item Large receptive fields
    \item Low spatial resolution
\end{itemize}
\end{block}

\vspace{0.2cm}

\begin{block}{Shallow layers}
\begin{itemize}
    \item High spatial resolution
    \item Precise localization information
    \item Weaker semantic representation
\end{itemize}
\end{block}

\end{column}

\begin{column}{0.47\textwidth}

\begin{block}{The trade-off}
\centering
\[
\text{Semantic information}
\quad \leftrightarrow \quad
\text{Spatial precision}
\]
\end{block}

\vspace{0.3cm}

\begin{itemize}
    \item Using only deep features:
    \textbf{good semantics, poor localization}
    
    \item Using only shallow features:
    \textbf{good localization, weaker semantics}
\end{itemize}

\end{column}

\end{columns}

\vspace{0.3cm}

\centering
\textbf{Object detection needs both.}

\end{frame}




%------------------------------------------------------------------------------

%------------------------------------------------------------------------------

%------------------------------------------------------------------------------

\begin{frame}{How the Feature Pyramid Network Works}

\begin{columns}[t]

\begin{column}{0.40\textwidth}

\begin{block}{Deep layers}
\footnotesize
\begin{itemize}
    \item Strong semantic information
    \item Large receptive fields
    \item Lower spatial resolution
\end{itemize}
\end{block}

\vspace{0.15cm}

\begin{block}{Shallow layers}
\footnotesize
\begin{itemize}
    \item Higher spatial resolution
    \item Better localization detail
    \item Weaker semantic representation
\end{itemize}
\end{block}

\vspace{0.15cm}

\centering
\textbf{FPN combines both}

\end{column}

\begin{column}{0.57\textwidth}

\centering

\placeholderimage[\textwidth]{Images/fpn_general.png}

\end{column}

\end{columns}

\vspace{0.15cm}

\centering
\textbf{Top-down pathway + lateral connections}

\end{frame}



%------------------------------------------------------------------------------

\begin{frame}{Reference Pneumonia Framework vs. FCOS}

\begin{columns}[t]

\begin{column}{0.48\textwidth}

\centering
\textbf{Reference pneumonia framework}

\vspace{0.2cm}

\placeholderimage[0.92\textwidth]{Images/reference_pneumonia_fpn.png}

\vspace{0.2cm}

{\footnotesize
Five backbone outputs are independently projected
via $1\times1$ convolutions.
}

\end{column}

\begin{column}{0.48\textwidth}

\centering
\textbf{FCOS}

\vspace{0.2cm}

\placeholderimage[0.6\textwidth]{Images/fcos_fpn.png}

\vspace{0.2cm}

{\footnotesize
$C_3$--$C_5$ are fused through a top-down FPN,
then extended with $P_6$ and $P_7$.
}

\end{column}

\end{columns}

\vspace{0.35cm}

\centering
\textbf{The project's implementation follows the FCOS feature construction.}

\end{frame}
%------------------------------------------------------------------------------
%------------------------------------------------------------------------------

\begin{frame}{FCOS-Style Detection Head}

\begin{columns}[t]

\begin{column}{0.43\textwidth}

\begin{itemize}

    \item One independent head for each
    \textbf{FPN level} $P_3$--$P_7$

    \item Each input feature map has
    \textbf{256 channels}

    \item Two feature towers:
    \begin{itemize}
        \item Classification
        \item Regression
    \end{itemize}

\end{itemize}

\vspace{0.2cm}

\centering
\textbf{Per-location predictions}

\[
\text{class}
\qquad
(l,t,r,b)
\qquad
\text{centerness}
\]

\end{column}

\begin{column}{0.55\textwidth}

\centering

\placeholderimage[1\textwidth]{Images/detection_head.png}

\end{column}

\end{columns}

\end{frame}


%------------------------------------------------------------------------------

\begin{frame}{What Does Each Head Predict?}

\begin{columns}[t]

\begin{column}{0.31\textwidth}

\begin{block}{Classification}
\centering
\[
1 \times H_i \times W_i
\]

\vspace{0.1cm}

\footnotesize
Is this location associated with pneumonia?
\end{block}

\end{column}

\begin{column}{0.31\textwidth}

\begin{block}{Box Regression}
\centering
\[
4 \times H_i \times W_i
\]

\vspace{0.1cm}

\footnotesize
$(l,t,r,b)$ distances
\end{block}

\end{column}

\begin{column}{0.31\textwidth}

\begin{block}{Centerness}
\centering
\[
1 \times H_i \times W_i
\]

\vspace{0.1cm}

\footnotesize
Quality of the spatial location
\end{block}

\end{column}

\end{columns}

\vspace{0.5cm}

\centering

\[
P_i
\quad\longrightarrow\quad
\boxed{\text{Classification}}
\quad+\quad
\boxed{\text{Regression}}
\quad+\quad
\boxed{\text{Centerness}}
\]

\vspace{0.25cm}

\begin{block}{Key point}
\centering
\textbf{Predictions are associated directly with spatial locations,
not anchor boxes.}
\end{block}

\end{frame}

%------------------------------------------------------------------------------

%------------------------------------------------------------------------------

%------------------------------------------------------------------------------

%------------------------------------------------------------------------------

\begin{frame}{Target Assignment}

\begin{columns}[t]

\begin{column}{0.47\textwidth}

\begin{block}{1. Map locations to the image}
\footnotesize

For location $(x,y)$ and stride $s$:

\[
p_x=(x+0.5)s,
\qquad
p_y=(y+0.5)s
\]

A location is positive only if it lies
inside a ground-truth box.

\end{block}

\vspace{0.12cm}

\begin{block}{2. Build LTRB targets}
\footnotesize

\[
l=p_x-x_1,\qquad t=p_y-y_1
\]

\[
r=x_2-p_x,\qquad b=y_2-p_y
\]

\end{block}

\end{column}

\begin{column}{0.47\textwidth}

\begin{block}{3. Assign the FPN level}
\footnotesize

\[
d_{\max}=\max(l,t,r,b)
\]

\[
\begin{array}{c|c}
P_3 &[0,64)\\
P_4 &[64,128)\\
P_5 &[128,256)\\
P_6 &[256,512)\\
P_7 &[512,\infty)
\end{array}
\]

\end{block}

\vspace{0.12cm}

\begin{block}{4. Centerness target}
\footnotesize

\[
c^*=
\sqrt{
\frac{\min(l,r)}{\max(l,r)}
\frac{\min(t,b)}{\max(t,b)}
}
\]

\end{block}

\end{column}

\end{columns}

\vspace{0.05cm}

\centering
{\scriptsize
\textbf{GT boxes} $\;\longrightarrow\;$ 
\textbf{positive mask + LTRB + centerness targets}
}

\end{frame}


%------------------------------------------------------------------------------

\begin{frame}{Loss Functions}

\begin{block}{Classification --- Sigmoid Focal Loss}

\[
\mathcal{L}_{\mathrm{cls}}
=
-\alpha_t(1-p_t)^\gamma \log(p_t)
\]

\centering
{\footnotesize $\alpha=0.25,\qquad \gamma=2$}

\end{block}

\vspace{0.15cm}

\begin{block}{Bounding Box --- GIoU Loss}

\[
\operatorname{GIoU}(A,B)
=
\operatorname{IoU}(A,B)
-
\frac{|C\setminus(A\cup B)|}{|C|}
\]

\[
\mathcal{L}_{\mathrm{reg}}=1-\operatorname{GIoU}(A,B)
\]

\end{block}

\vspace{0.15cm}

\begin{block}{Centerness --- Sigmoid Focal Loss}

\[
\mathcal{L}_{\mathrm{ctr}}
=
-\alpha_t(1-p_t)^\gamma\log(p_t)
\]

\end{block}

\vspace{0.12cm}

\centering
\[
\boxed{
\mathcal{L}
=
\mathcal{L}_{\mathrm{cls}}
+
\mathcal{L}_{\mathrm{reg}}
+
\mathcal{L}_{\mathrm{ctr}}
}
\]

{\footnotesize Equal weights for the three components.}

\end{frame}

%==============================================================================
% 3. EXPERIMENTAL SETUP
%==============================================================================

\section{Experimental Setup}

%------------------------------------------------------------------------------

\begin{frame}{Experimental Setup}

\begin{columns}[t]

\begin{column}{0.48\textwidth}

\begin{block}{Dataset}
\small
\begin{itemize}
    \item RSNA Pneumonia Detection Challenge
    \item Bounding-box annotations
    \item Input resolution: $512\times512$
\end{itemize}
\end{block}

\vspace{0.2cm}

\begin{block}{Common Pipeline}
\small
\begin{itemize}
    \item Same detector architecture
    \item Same data split
    \item Same target assignment
    \item Same evaluation procedure
\end{itemize}
\end{block}

\end{column}

\begin{column}{0.48\textwidth}

\begin{block}{Backbone Configurations}
\centering
\begin{tabular}{lcc}
\toprule
 & \textbf{ImageNet} & \textbf{Chest-Xray} \\
\midrule
ResNet-50  & R50-IN  & R50-CX \\
ResNet-101 & R101-IN & R101-CX \\
\bottomrule
\end{tabular}
\end{block}

\vspace{0.25cm}

\centering
\[
\boxed{
2\ \text{depths}
\times
2\ \text{initializations}
=
4\ \text{final models}
}
\]

\end{column}

\end{columns}

\vspace{0.3cm}



\centering
\textbf{Goal: study the effect of backbone depth and initialization.}

\end{frame}

%------------------------------------------------------------------------------
\begin{frame}{Training Strategy}
\begin{columns}[t]

\begin{column}{0.46\textwidth}

\begin{block}{Optimization}
\small
\begin{itemize}
    \item Backbone freezing
    \item Differential learning rates
    \item Learning-rate warm-up
    \item Cosine annealing
\end{itemize}
\end{block}

\vspace{0.25cm}

\begin{block}{Stability}
\small
\begin{itemize}
    \item Gradient clipping
    \item Exponential Moving Average (EMA)
    \item Validation-based checkpoint selection
\end{itemize}
\end{block}

\end{column}

\begin{column}{0.46\textwidth}
\centering

\[
\text{Freeze}
\rightarrow
\text{Warm-up}
\rightarrow
\text{Unfreeze}
\rightarrow
\text{Fine-tune}
\]

\vspace{0.45cm}

\begin{alertblock}{Main idea}
\centering
\small
Protect pretrained features first,\\
then progressively adapt the detector.
\end{alertblock}

\vspace{0.3cm}

{\footnotesize
Training schedules were adapted to the convergence behaviour
of each configuration.
}

\end{column}

\end{columns}
\end{frame}







\begin{frame}{Why Longer Training?}

\centering

\vspace{0.15cm}

\placeholderimage[0.45\textwidth]{Images/AP/why_long.png}

\vspace{0.2cm}

\begin{alertblock}{ResNet-101}
\centering
\small
Slower convergence and larger validation fluctuations
\end{alertblock}

\vspace{0.1cm}

{\footnotesize
Simply increasing the number of epochs was not sufficient.
}
\end{frame}

































\begin{frame}{Evaluation Metrics}

\begin{columns}[t]

\begin{column}{0.31\textwidth}

\begin{block}{Detection}
\small
\textbf{AP}\\
Average Precision across confidence thresholds

\vspace{0.15cm}

\textbf{AR@10}\\
Average Recall with at most 10 detections/image
\end{block}

\end{column}

\begin{column}{0.31\textwidth}

\begin{block}{Image-level}
\small
\textbf{Precision}\\
Fraction of positive predictions that are correct

\vspace{0.1cm}

\textbf{Recall}\\
Fraction of positive images correctly detected

\vspace{0.1cm}

\textbf{F1-score}\\
Balance between precision and recall
\end{block}

\end{column}

\begin{column}{0.31\textwidth}

\begin{block}{Localization}
\small
\textbf{Mean IoU}\\
Average overlap of matched boxes

\vspace{0.1cm}

\textbf{Box F1}\\
Box-level precision--recall balance
\end{block}

\end{column}

\end{columns}

\vspace{0.4cm}

\begin{alertblock}{Operating point}
\centering
\small
For image-level metrics, the confidence threshold is selected on the
validation set by maximizing Youden's $J$.
\end{alertblock}

\vspace{0.1cm}

\centering
\[
J = \mathrm{Sensitivity}+\mathrm{Specificity}-1
\]

\end{frame}
%==============================================================================
% 4. RESULTS
%==============================================================================

\section{Results}

%------------------------------------------------------------------------------

\begin{frame}{Quantitative Results}

\begin{center}
\small

\begin{tabular}{lccccc}
\toprule
\textbf{Model} &
\textbf{AP} &
\textbf{AR@10} &
\textbf{Image F1} &
\textbf{Mean IoU} &
\textbf{Box F1} \\
\midrule
R50-IN  & \textbf{0.423} & 0.886 & 0.654 & 0.688 & \textbf{0.448} \\
R50-CX  & 0.399 & 0.876 & 0.633 & 0.684 & 0.424 \\
R101-IN & 0.397 & \textbf{0.901} & \textbf{0.668} &
          \textbf{0.692} & 0.446 \\
R101-CX & 0.239 & 0.804 & 0.649 & 0.659 & 0.388 \\
\bottomrule
\end{tabular}

\end{center}

\vspace{0.45cm}

\begin{columns}

\begin{column}{0.46\textwidth}

\begin{block}{R50-IN}
\centering
Highest AP and slightly highest Box F1.
\end{block}

\end{column}

\begin{column}{0.46\textwidth}

\begin{block}{R101-IN}
\centering
Highest AR@10, Image F1 and mean matched IoU.
\end{block}

\end{column}

\end{columns}

\end{frame}

\begin{frame}{Precision--Recall and Object Size}

\begin{columns}[t]

\begin{column}{0.56\textwidth}
\centering

\placeholderimage[0.63\linewidth]{Images/comparison/precision_recall_comparison.png}

\vspace{0.08cm}

{\scriptsize
Precision--Recall curves on the validation set
}

\vspace{0.18cm}

\begin{alertblock}{Key observation}
\centering
\small
R50-IN provides the strongest integrated
Precision--Recall trade-off.
\end{alertblock}

\end{column}

\begin{column}{0.40\textwidth}

\begin{block}{Overall AP}
\centering
\small
\begin{tabular}{lc}
\toprule
\textbf{Model} & \textbf{AP} \\
\midrule
R50-IN  & \textbf{0.423} \\
R50-CX  & 0.399 \\
R101-IN & 0.397 \\
R101-CX & 0.239 \\
\bottomrule
\end{tabular}
\end{block}

\vspace{0.18cm}

\begin{block}{By region size}
\centering
\scriptsize
\begin{tabular}{lcc}
\toprule
\textbf{Model} & $\mathbf{AP_M}$ & $\mathbf{AP_L}$ \\
\midrule
R50-IN  & 0.039 & \textbf{0.442} \\
R50-CX  & 0.024 & 0.433 \\
R101-IN & 0.035 & 0.413 \\
R101-CX & 0.029 & 0.231 \\
\bottomrule
\end{tabular}
\end{block}

\vspace{0.18cm}

\begin{alertblock}{Scale-related limitation}
\centering
\small
$\mathbf{AP_M < 0.04}$ for all four models.
\\[0.05cm]
Medium-sized regions remain difficult.
\end{alertblock}

\end{column}

\end{columns}

\end{frame}
\begin{frame}{Image-Level Results: ImageNet Initialization}

\begin{columns}[c]

\begin{column}{0.48\textwidth}
\centering

\textbf{R50-IN}

\vspace{0.15cm}

\placeholderimage[0.7\textwidth]{Images/confusion/R50-IN.png}

\vspace{0.05cm}

{\small
F1 = 0.654 \qquad Youden $J$ = 0.628
}

\end{column}

\begin{column}{0.48\textwidth}
\centering

\textbf{R101-IN}

\vspace{0.15cm}

\placeholderimage[0.7\textwidth]{Images/confusion/R101-IN.png}

\vspace{0.05cm}

{\small
F1 = \textbf{0.668} \qquad Youden $J$ = \textbf{0.646}
}

\end{column}

\end{columns}

\vspace{0.25cm}

\begin{alertblock}{Observation}
\centering
\small
R101-IN provides the stronger image-level operating point.
\end{alertblock}

\end{frame}
\begin{frame}{Image-Level Results: Chest-Xray Initialization}

\begin{columns}[c]

\begin{column}{0.48\textwidth}
\centering

\textbf{R50-CX}

\vspace{0.15cm}

\placeholderimage[0.7\textwidth]{Images/confusion/R50-CX.png}

\vspace{0.05cm}

{\small
F1 = 0.633 \qquad Youden $J$ = 0.602
}

\end{column}

\begin{column}{0.48\textwidth}
\centering

\textbf{R101-CX}

\vspace{0.15cm}

\placeholderimage[0.7\textwidth]{Images/confusion/R101-CX.png}

\vspace{0.05cm}

{\small
F1 = 0.649 \qquad Youden $J$ = 0.616
}

\end{column}

\end{columns}

\vspace{0.25cm}

\begin{alertblock}{Observation}
\centering
\small
ImageNet initialization performs better than Chest-Xray initialization
under the final training conditions.
\end{alertblock}

\end{frame}


\begin{frame}{Qualitative Example: R101-IN}

\begin{columns}[c]

\begin{column}{0.46\textwidth}
\centering

\textbf{Input radiograph}

\vspace{0.15cm}

\placeholderimage[0.6\textwidth]{Images/qualitative/r101in_input.jpg}

\end{column}

\begin{column}{0.46\textwidth}
\centering

\textbf{Feature flow and prediction}

\vspace{0.15cm}

\placeholderimage[1\textwidth]{Images/qualitative/r101in_feature_flow.png}

\end{column}

\end{columns}

\vspace{0.25cm}

\begin{alertblock}{Qualitative inspection}
\centering
\small
The detector progressively transforms the radiograph into
multi-scale feature representations before producing the final prediction.
\end{alertblock}

\end{frame}




\begin{frame}{Clinical Review: Limitations of the Dataset}

\begin{columns}[t]

\begin{column}{0.48\textwidth}

\begin{block}{Image acquisition}
\small
\begin{itemize}
    \item Many examinations are \textbf{portable chest X-rays}
    \item Patients are often examined in a \textbf{supine or non-standing position}
    \item Positioning can change the apparent extent and distribution of opacities
    \item Only \textbf{one projection} is available in the dataset
\end{itemize}
\end{block}

\end{column}

\begin{column}{0.48\textwidth}

\begin{block}{Clinical assessment}
\small
\begin{itemize}
    \item In practice, radiologists can use \textbf{frontal and lateral projections}
    \item Diagnosis is based on the \textbf{overall thoracic appearance}, not only on the opacity
    \item Cardiac size, diaphragm and costophrenic angles can provide additional information
    \item Clinical history and previous imaging may further change the interpretation
\end{itemize}
\end{block}

\end{column}

\end{columns}

\vspace{0.3cm}

\begin{alertblock}{Main limitation}
\centering
\small
The model receives \textbf{a single radiograph and an opacity annotation},
while clinical assessment can integrate \textbf{multiple projections and additional anatomical and clinical information}.
\end{alertblock}

\end{frame}





\begin{frame}{Implementation, Hardware & Resources}

\begin{columns}[t]

\begin{column}{0.48\textwidth}

\begin{block}{Software}
\small
\begin{itemize}
    \item PyTorch
    \item torchvision
\end{itemize}
\end{block}

\vspace{0.2cm}

\begin{block}{Datasets}
\small
\textbf{RSNA Pneumonia Detection Challenge}\\
Used for detector training and evaluation.

\vspace{0.12cm}

\textbf{Chest X-Ray Images (Pneumonia)}\\
Used for Chest-Xray pretraining.
\end{block}

\end{column}

\begin{column}{0.48\textwidth}

\begin{block}{Hardware}
\small
\begin{itemize}
    \item NVIDIA RTX 4090
    \item Cloud GPU rental (\url{https://www.runpod.io/})
    \item $\sim$20 experiments
    \item Total cost: $\sim$\textbf{\euro20}
\end{itemize}
\end{block}

\vspace{0.2cm}

\begin{block}{Repository}
\centering
\small
\url{https://github.com/Francesco2035/}\\
\url{fcos-pneumonia-detection-naml-25-26}
\end{block}

\end{column}

\end{columns}

\end{frame}




\begin{frame}{Conclusions}

\begin{block}{1. Anchor-free detection}
\centering
A fully convolutional FCOS-style detector can localize pneumonia
without predefined anchor boxes.
\end{block}

\vspace{0.2cm}

\begin{block}{2. Initialization and backbone}
\centering
ImageNet initialization provided the most reliable starting point.
The Chest-Xray pretrained models were relatively disadvantageous,
especially for ResNet-101.
\end{block}

\vspace{0.2cm}

\begin{block}{3. Main lesson}
\centering
A strong backbone is important, but \textbf{careful training and optimization}
have a major impact on the final performance.
Better-controlled training and richer clinical information remain key directions.
\end{block}

\vspace{0.3cm}

\centering
\textbf{Best overall operating-point model: R101-IN}

\end{frame}




\begin{frame}{References}

\footnotesize

\begin{thebibliography}{9}

\bibitem{wu2024}
L. Wu et al.,
\textit{Pneumonia detection based on RSNA dataset and anchor-free
deep learning detector},
Scientific Reports, 2024.

\bibitem{fcos}
Z. Tian, C. Shen, H. Chen, T. He,
\textit{FCOS: Fully Convolutional One-Stage Object Detection},
ICCV, 2019.

\bibitem{resnet}
K. He, X. Zhang, S. Ren, J. Sun,
\textit{Deep Residual Learning for Image Recognition},
CVPR, 2016.

\bibitem{fpn}
T.-Y. Lin et al.,
\textit{Feature Pyramid Networks for Object Detection},
CVPR, 2017.

\bibitem{focal}
T.-Y. Lin et al.,
\textit{Focal Loss for Dense Object Detection},
ICCV, 2017.

\bibitem{giou}
H. Rezatofighi et al.,
\textit{Generalized Intersection over Union: A Metric and a Loss
for Bounding Box Regression},
CVPR, 2019.

\end{thebibliography}

\end{frame}




%==============================================================================

\end{document}